\chapter{Results and Discussion}

\section{SpamAssassin}

SpamAssassin was run three times over the 681-email corpus: a no-delay
baseline, Run~2 (delay, default timeouts), and Run~3 (delay, reduced timeouts).
Each message was scanned once by the standalone binary and timed from arrival to
exit.

\begin{table}[h]\centering
\small
\begin{tabular}{@{}lccc@{}}
\toprule
Metric & Baseline & Run 2 & Run 3 \\
\midrule
Flagged YES           & 383 (56.2\%) & 355 (52.1\%) & 316 (46.4\%) \\
Emails with DNSBL hit  & 416 & 400 & 251 \\
Scan min/avg/max (s)   & 0.45/1.68/7.58 & 3.17/7.36/16.50 & 0.43/4.05/14.83 \\
\bottomrule
\end{tabular}
\caption{SpamAssassin - Aggregate results across the three runs.}
\label{tab:sa-agg}
\end{table}

\subsection{Detection}
\label{sec:spamassassin_detection}

With no delay, 383 messages were flagged. Applying 1.3\,s each way then broke
detection at the resolver: DNSBL answers returned at about 2.64\,s (one RTT). But,
they were discarded by \texttt{unbound} before reaching the scanner with the error message \texttt{drop reply, it is older than discard-timeout}. The default value of \texttt{discard-timeout}, 1900\,ms, is lower than the RTT of 2600\,ms, thus all delayed responses were too old and discarded~\cite{archwiki_unbound}. Setting \texttt{discard-timeout} to 9000 and \texttt{infra-cache-min-rtt}
to 4000, which were higher than the RTT, solved the problem, and a cold
resolution took 2.694\,s while the test message regained its score.

With the resolver fixed, Run~2 nearly preserved detection (416\,$\to$\,400 DNSBL
emails) at about four times the scan time. Run~3 cut \texttt{rbl\_timeout} and
\texttt{spf\_timeout} to 1.5\,s, below the RTT, trading detection for speed. Two
things follow. First, every DNSBL rule fired less than in Run~2 and none fired
more (Table~\ref{tab:sa-rules}).

\begin{table}[H]\centering
\small
\begin{tabular}{@{}lccc@{}}
\toprule
Rule & Baseline & Run 2 & Run 3 \\
\midrule
RCVD\_IN\_PBL          & 151 & 149 & 73 \\
RCVD\_IN\_SBL\_CSS     & 148 & 136 & 43 \\
URIBL\_DBL\_SPAM       & 118 & 117 & 73 \\
RCVD\_IN\_PSBL         & 108 & 93  & 67 \\
RCVD\_IN\_MSPIKE\_BL   & 103 & 93  & 32 \\
RCVD\_IN\_XBL          & 82  & 75  & 20 \\
RCVD\_IN\_BL\_SPAMCOP\_NET & 75 & 63 & 12 \\
\bottomrule
\end{tabular}
\caption{SpamAssassin - DNSBL/URIBL rule firing (selected rules); every rule falls from Run~2
to Run~3.}
\label{tab:sa-rules}
\end{table}

Each rule corresponds to one blocklist: RCVD\_IN\_* rules check the sending server's IP against an address-based list, and URIBL\_* rules check domains in the message body against a domain-based list. (Table~\ref{tab:sa-rules}) shows the individual DNSBL/URIBL rules; every rule falls from Run 2 to Run 3.

Second, the resolver failed after approximately email~500. With
\texttt{rbl\_timeout} and \texttt{spf\_timeout} reduced to 1.5\,s,
SpamAssassin abandoned DNSBL lookups before the replies could return across the
2.6\,s DTN round-trip. However, the delayed replies still reached
\texttt{unbound}. From the resolver's perspective, the upstream DNS server had
failed to respond before its retransmission timeout (RTO), so the query was
treated as an upstream timeout and the RTO for that server was increased using
exponential backoff. As this occurred repeatedly across hundreds of emails, the
RTO eventually reached Unbound's maximum value of 120\,000\,ms, as defined by
the default \texttt{infra-cache-max-rtt} setting~\cite{archwiki_unbound}.

According to the Unbound documentation~\cite{unbound_timeout}, reaching this value places the upstream server into the \emph{blocking regime},  where Unbound considers the upstream server unresponsive and temporarily stops sending DNS queries to it and Unbound will attempt to use the server again only after the infrastructure-cache entry expires. Because the experiment used a single upstream resolver, subsequent DNSBL lookups could no longer be resolved successfully. Consequently, after approximately email~500, SpamAssassin effectively scanned messages without DNS-based checks, causing the DNSBL hit rate to collapse (Table~\ref{tab:sa-collapse}) and scan times to fall sharply. The resolver remained in this blocked state until it was restarted.

\begin{table}[H]\centering
\small
\begin{tabular}{@{}lcc@{}}
\toprule
Corpus position & Run 2 & Run 3 \\
\midrule
Emails 1--500   & 260/500 (52\%) & 220/500 (44\%) \\
Emails 501--681 & 140/181 (77\%) & 31/181 (17\%) \\
\bottomrule
\end{tabular}
\caption{SpamAssassin - DNSBL-hit rate by corpus position; Run~3 collapses after
$\approx$email 500.}
\label{tab:sa-collapse}
\end{table}

The 251 DNSBL-hit emails in Run 3, despite a 1.5 s timeout below the 2.6 s RTT, are explained by the adaptive nature of rbl\_timeout. It is not a hard per-query cutoff but the starting value of a timeout that shrinks as queries complete, with a one-second grace floor granted to each still-pending query. Messages whose lookups form a serial dependency chain — the URIBL domain→NS→IP→DNSBL wave chain, or nested SPF include: directives — launch fresh pending queries at each new wave, and the grace floor keeps extending the effective deadline past the RTT, so the delayed answers arrive in time to score. The correlation confirms this: emails that scanned longer than the RTT hit a DNSBL rule 47\% of the time, against 24\% for emails that scanned faster, and 83\% of all Run 3 hits came from the longer-scanning group. Cache reuse of shared-domain listings is a secondary contributor. The result is that the same serial dependency chains that increase scan time under network delay also preserve DNSBL detection when the timeout is reduced, because the longer scan duration allows delayed DNS replies to arrive before scanning completes.

\subsection{Scan time}

Even in baseline conditions, the scan is not instantaneous, the network rules depend
on the DNSBL replies, which can take time, and the scan concludes only when the
last one of them arrives; with the resolver set up properly, delayed scans were still taking around 8 seconds. Since the lookups go in serial waves, dependent on the results from the previous waves, this suggests that the latency depends on the depth of the dependency chain, not the number of lookups. There are two causes of such depth – URIBL chain, which resolves the name servers of the domain into IPs and performs DNSBL lookup of those IPs (domain > NS > IP > DNSBL), and nested SPF \texttt{include:} directives, each fetched before its children are visible.

\begin{table}[H]\centering
\small
\begin{tabular}{@{}lccccc@{}}
\toprule
Run & $<$1\,s & 1--3\,s & 3--6\,s & 6--10\,s & $>$10\,s \\
\midrule
Baseline & 438 & 101 & 104 & 38  & 0 \\
Run 2    & 0   & 0   & 316 & 213 & 152 \\
Run 3    & 162 & 129 & 194 & 184 & 12 \\
\bottomrule
\end{tabular}
\caption{SpamAssassin - Scan-time distribution (emails per bucket).}
\label{tab:sa-timing}
\end{table}

The distribution changes completely with delay, since no Run~2 message scan completes in less than 3\,s, since all must now do at least one round-trip time and dependency-chain messages must do many. Run 3 splits into two groups: the shorter timeout cuts many scans short, while dependency-chain messages still run long.

\subsection{Flips}
\label{sec:sa-flips}

\begin{table}[H]\centering
\small
\begin{tabular}{@{}lc@{}}
\toprule
 & Run 2 $\to$ Run 3 \\
\midrule
Verdict unchanged           & 424 \\
Flipped YES $\to$ NO        & 148 \\
Flipped NO $\to$ YES        & 109 \\
Net flagged change          & $-39$ \\
Gained a DNSBL rule (drift) & 178 (26.1\%) \\
\bottomrule
\end{tabular}
\caption{SpamAssassin - Verdict changes and blocklist drift, Run~2 to Run~3.}
\label{tab:sa-flips}
\end{table}

Run~3 was done two days later than Run~2, so the effect for this pair is a combination of both. The shortened timeout aborts the DNS query before the answer comes back, eliminating DNSBL matches and shifting messages from YES to NO. However, 178 messages (26.1\%) \emph{acquired} the DNSBL rule between the runs, and a shortened timeout cannot do this, so this is blocklist drift over the two days. The 109 NO\,$\to$\,YES switches are thus drift, while the 148 YES\,$\to$\,NO switches are a mixture of timeout failure and drift. 

\subsection{Network volume}

\begin{table}[H]\centering
\small
\begin{tabular}{@{}lccc@{}}
\toprule
Component & Baseline & Run 2 & Run 3 \\
\midrule
DNS bytes       & 5{,}160{,}681 & 5{,}418{,}742 & 3{,}413{,}349 \\
DNS packets     & 37{,}789 & 39{,}875 & 25{,}783 \\
Email Bundle bytes    & 572{,}625 & 538{,}529 & 547{,}709 \\
Email Bundle packets  & 1{,}447 & 1{,}369 & 1{,}389 \\
DNS \% of bytes & 90.0\% & 91.0\% & 86.2\% \\
\bottomrule
\end{tabular}
\caption{SpamAssassin - Traffic on the relay: DNS to the resolver vs DTN bundle traffic.}
\label{tab:sa-bw}
\end{table}

Traffic in DNS is nine to ten times greater than the DTN bundle traffic carrying
the email messages in the baseline and Run~2, while the amount of
email bundle traffic stays almost the same across all runs. In Run~3, the lower
DNS traffic indicates a resolver failure rather than a bandwidth saving. After
about email~500, the resolver stopped sending DNS queries upstream, so fewer
DNS packets were generated and SpamAssassin could no longer perform DNSBL
lookups.

\section{rspamd}

This section reports the three runs with rspamd. The runs follow the same corpus and DNS path used by the previous tool. rspamd was executed on the
receiver using \texttt{rspamc --ip 1.2.3.4 --mime}; the \texttt{--ip} option ensures that the external processing is done, and the rules dependent on DNS run. Runs~1
and~2 were executed using the default task timeout, which was 8\,s,
while Run~3 uses 1.5\,s, below the round-trip time.

\subsection{Detection}

\begin{table}[H]\centering
\small
\begin{tabular}{@{}lccc@{}}
\toprule
Metric & Run 1 (baseline) & Run 2 (delay) & Run 3 (delay, 1.5\,s) \\
\midrule
Messages processed            & 681 & 681 & 681 \\
Flagged as spam               & 675 (99.1\%) & 654 (96.0\%) & 569 (83.6\%) \\
Messages with real hit        & 502 & 482 & 94 \\
Messages with \texttt{\_FAIL} & 0 & 45 & 665 \\
Total real hits               & 1449 & 1330 & 221 \\
Total \texttt{\_FAIL}         & 0 & 357 & 2729 \\
Average score                 & 18.13 & 17.20 & 9.26 \\
Maximum score                 & 46.10 & 45.55 & 21.55 \\
\bottomrule
\end{tabular}
\caption{rspamd - Detection and scoring across the three runs.}
\label{tab:rspamd-detection}
\end{table}

Baseline flags 675 of 681 messages (99.1\%); this is the detection ceiling.
Under delay with default timeouts (Run~2) detection is nearly preserved at 654
(96.0\%). Reducing \texttt{task\_timeout} to 1.5\,s (Run~3) drops it to 569
(83.6\%). The cause is in the hit and \texttt{\_FAIL} rows: real hits fall
502\,$\to$\,482\,$\to$\,94 as timed-out lookups rise 0\,$\to$\,45\,$\to$\,665.


\subsection{Scan time}

\begin{table}[H]\centering
\small
\begin{tabular}{@{}lccc@{}}
\toprule
Scan time (s) & Run 1 & Run 2 & Run 3 \\
\midrule
$<$1\,s   & 644 & 0   & 2 \\
1--2\,s   & 24  & 0   & 679 \\
2--4\,s   & 1   & 113 & 0 \\
4--6\,s   & 6   & 428 & 0 \\
$\geq$6\,s & 6  & 140 & 0 \\
\bottomrule
\end{tabular}
\caption{rspamd - Scan-time distribution across the three runs (messages per bucket).}
\label{tab:rspamd-scan}
\end{table}

Baseline scans fast: 644 out of 681 scans in less than 1\,s. In the delay setup with
default timeouts (Run~2) every scan is slow (median 5.67\,s), where the round trip time
dominates since every uncached lookup incurs one RTT cost; 140 out of 681 scans
hit the ceiling at 8\,s. In Run~3, scan time reduced to 1 to
2\,s for 679 out of 681 scans, with none above 1.64\,s.

\subsection{Flips, Run 2 to Run 3}

\begin{table}[H]\centering
\small
\begin{tabular}{@{}lc@{}}
\toprule
 & Run 2 $\to$ Run 3 \\
\midrule
Verdict unchanged        & 560 \\
Flipped YES $\to$ NO     & 103 \\
Flipped NO $\to$ YES     & 18 \\
Net flagged change       & $-85$ \\
Messages that lost score & 579 \\
Mean score change        & $-7.94$ \\
\bottomrule
\end{tabular}
\caption{rspamd - Per-message flips between the two delayed runs.}
\label{tab:rspamd-flips}
\end{table} 
Reducing \texttt{task\_timeout} flipped 103 messages from spam to not-spam and
lowered the score of 579 of 681. The 18 messages that gained detection run
counter to a shorter timeout and are most likely blocklist-listing drift between
the two runs rather than a timeout effect.

\subsection{Network volume}

\begin{table}[H]\centering
\small
\begin{tabular}{@{}lccc@{}}
\toprule
Metric (receiver capture) & Run 1 & Run 2 & Run 3 \\
\midrule
DNS queries sent          & 38{,}873 & 80{,}890 & 24{,}116 \\
DNS responses received    & 38{,}771 & 80{,}822 & 24{,}078 \\
Queries to upstream       & 13{,}894 & 19{,}033 & 5{,}969 \\
Total packets             & 89\,k & 182\,k & 59\,k \\
Total data                & 18\,MB & 32\,MB & 14\,MB \\
\bottomrule
\end{tabular}
\caption{rspamd - Network volume per run.}
\label{tab:rspamd-network}
\end{table}

The same number of sent and received queries is seen for all runs, meaning that there were no lost queries on the wire, while changes in traffic volume are due to rspamd itself, not to packet loss.

Run~2 sent almost twice the number of queries than the baseline (38{,}873\,$\to$\,80{,}890), and there was a proportional increase in data size (18\,$\to$\,32\,MB). The reason is retransmission. rspamd retries sending the DNS request every time it times out waiting for the response, with the number of retransmits equal to the retransmit count, which gives a total per-request timeout of \texttt{timeout\,$\times$\,retransmits}, which is 1\,s $\cdot$ 5 by default. As the timeout per attempt is shorter than the 2.6\,s round trip time, the unanswered queries will be resent until they get their response. So, an entirely fresh lookup will be sent 3 times on average until the answer is received at 2.6\,s.


Run~3 does the opposite: query volume falls below the baseline (80{,}890$\,\rightarrow$\,24{,}116 queries, 32$\,\rightarrow$\,14\,MB). Since the timeout is only 1.5 seconds, the task is terminated before the DNS budget expires (it is 5 seconds). Therefore, a request can be issued only twice before termination, stopping the sequence of retransmissions that inflated
Run~2. This lower volume is the network-side view of the 2{,}729 \texttt{\_FAIL} symbols: the queries went out but were abandoned before
their answers returned.


\section{Vipul's Razor}

This section reports the three runs with Razor. The runs follow the same corpus
and DNS path used by the previous tools. Razor 2.84 was executed on the receiver
using \texttt{razor-check < message.eml}. Razor exposes no network timeout through
\texttt{razor-agent.conf}, so the connect timeout is a hardcoded constant in
\texttt{Razor2::Client::Core}. Runs~1 and~2 used the default of 20\,s, while
Run~3 patched it to 2\,s, below the round-trip time.

\subsection{Detection}

\begin{table}[H]\centering
\small
\begin{tabular}{@{}lccc@{}}
\toprule
Metric & Run 1 (baseline) & Run 2 (delay) & Run 3 (delay, patched) \\
\midrule
Messages processed        & 681 & 681 & 681 \\
YES (exit 0, spam)        & 22 (3.23\%)  & 19 (2.79\%)  & 0 (0\%) \\
NO (exit 1, not spam)     & 659 (96.77\%) & 656 (96.33\%) & 0 (0\%) \\
FAIL (exit 2, error)      & 0 (0\%)   & 6 (0.88\%)   & 681 (100\%) \\
\bottomrule
\end{tabular}
\caption{Razor - Verdict distribution across the three runs.}
\label{tab:razor-detection}
\end{table}


Under delay with default timeouts (Run~2), the verdict remains essentially the same (3.23\%\,$\to$\,2.79\% flagged). Run~3 collapses the entire detection process
as it reports exit~2 in every scan performed. In contrast to the DNSBL-based tools, Razor has only one cooperative detection method, which means that if the connection fails, it reports no verdict at all.

\subsection{Scan time}

\begin{table}[H]\centering
\small
\begin{tabular}{@{}lccc@{}}
\toprule
Scan time (s) & Run 1 & Run 2 & Run 3 \\
\midrule
\multicolumn{4}{@{}l}{\textit{distribution (messages)}}\\
$<$1\,s      & 676 & 0   & 0 \\
1--5\,s      & 4   & 0   & 0 \\
5--10\,s     & 0   & 90  & 671 \\
10--15\,s    & 0   & 465 & 10 \\
15--30\,s    & 1   & 73  & 0 \\
30--60\,s    & 0   & 44  & 0 \\
$\geq$60\,s  & 0   & 9   & 0 \\
\bottomrule
\end{tabular}
\caption{Razor - Scan-time distribution across the three runs.}
\label{tab:razor-scan}
\end{table}


Baseline scans are completed within one second (676 of 681). With delay without timeout
(Run~2), the median increases to 11\,s, and the distribution gets a heavy tail up
to 122.6\,s. With the timeout (Run~3), the time of scanning becomes deterministic
with a value close to 8\,s, which means four connection attempts of the server,
each taking 2\,s.

\subsection{Flips}
\begin{table}[H]\centering
\small
\begin{tabular}{@{}lc@{}}
\toprule
Transition (Run 1 $\to$ Run 2) & Messages \\
\midrule
NO $\to$ NO (unchanged)   & 653 \\
YES $\to$ YES (unchanged) & 19 \\
NO $\to$ FAIL             & 6 \\
YES $\to$ NO              & 3 \\
\bottomrule
\end{tabular}
\hfill
\begin{tabular}{@{}lc@{}}
\toprule
Transition (Run 2 $\to$ Run 3) & Messages \\
\midrule
NO $\to$ FAIL   & 656 \\
YES $\to$ FAIL  & 19 \\
FAIL $\to$ FAIL & 6 \\
\bottomrule
\end{tabular}
\caption{Razor - Per-message flips, matched on a Message-ID and Delivery-date key.}
\label{tab:razor-flips}
\end{table}

Under delay alone the verdict is stable: 672 of 681 messages keep their baseline verdict, with 3 flipping YES→NO and 6 returning a network error. The timeout run flips 675 of 681 from a valid verdict to FAIL.


\subsection{Network volume}

\begin{table}[H]\centering
\small
\begin{tabular}{@{}lccc@{}}
\toprule
Metric (receiver capture) & Run 1 & Run 2 & Run 3 \\
\midrule
Frames captured              & 12{,}814 & 25{,}997 & 57{,}937 \\
Bytes captured               & 3.4\,MB  & 5.1\,MB  & 40.2\,MB \\
TCP connections to :2703     & 687      & 887      & 5{,}460 \\
Connections per message      & 1.01     & 1.30     & 8.02 \\
DNS queries / responses      & 36 / 42  & 650 / 810 & 777 / 971 \\
\bottomrule
\end{tabular}
\caption{Razor - Network volume per run.}
\label{tab:razor-network}
\end{table}


DNS traffic is less than 1\% of volume in all runs: Razor's per-check DNS traffic is
small, and its delay sensitivity is in TCP, not DNS. That is the opposite of rspamd, where DNS traffic dominated the delayed run.

Connections per message increase from 1.01\,$\to$\,1.30\,$\to$\,8.02. In the baseline, each message establishes, on average, one TCP connection to a Razor server, indicating that the first selected server usually responds successfully. With delay, some scans use additional servers due to failure; with the timeout, every handshake fails, so nextserver walks through every server in the catalogue and nomination lists until they are exhausted, and all of them are tried on each message. Therefore, the volume of Run~3 (57{,}937 frames, 40.2\,MB) is the largest despite the shortest wall-time: every handshake failure has a kernel-level SYN retransmission, which accounts for roughly two SYNs per attempt across four servers in 681 messages.


\section{Summary}

\begin{table}[H]\centering
\scriptsize
\renewcommand{\arraystretch}{1.3}
\begin{tabular}{@{}p{1.5cm} p{2.1cm} p{4.0cm} p{4.6cm}@{}}
\toprule
\textbf{Tool} & \textbf{Run} & \textbf{Observed} & \textbf{Reason} \\
\midrule

SpamAssassin & Run 2 (delay, default) &
Detection nearly preserved; scan time about four times higher &
Tuned resolver timers above the RTT let delayed DNSBL answers arrive; scans block waiting one RTT per lookup wave \\
\cmidrule(l){2-4}
& Run 3 (delay, reduced) &
Per-email detection falls; resolver collapses near email 500 &
Timeout below RTT abandons queries early; late answers drive unbound's rto to its ceiling, which does not self-recover \\
\midrule

rspamd & Run 2 (delay, default) &
Detection nearly preserved; DNS traffic roughly doubles &
Answers arrive within the 8\,s ceiling; 1\,s per-attempt DNS timeout below the RTT retransmits each uncached query \\
\cmidrule(l){2-4}
& Run 3 (delay, 1.5\,s) &
Detection drops; hits replaced by \texttt{\_FAIL}; DNS volume falls &
Task aborts before answers return, so lookups time out instead of listing; retransmit sequence is cut short \\
\midrule

Razor & Run 2 (delay, default) &
Detection essentially unchanged; scan time rises with a long tail &
Verdicts still returned over the link; slow servers hit the hardcoded read timeout, adding failover delay \\
\cmidrule(l){2-4}
& Run 3 (delay, patched) &
Detection collapses completely; every check returns FAIL &
2\,s connect timeout is below the 2.6\,s RTT, so every handshake aborts before its SYN-ACK returns \\
\midrule

All tools & Delay &
Filter DNS/TCP traffic dominates the link, far above the mail &
Network rules issue many lookups per message; the mail payload is small by comparison \\

\bottomrule
\end{tabular}
\caption{Summary of tool behaviour under delay.}
\label{tab:results-summary}
\end{table}






\begin{table}[H]\centering
\footnotesize
\renewcommand{\arraystretch}{1.4}
\begin{tabular}{@{}p{2.1cm} p{4.4cm} p{5.7cm}@{}}
\toprule
\textbf{Tool} & \textbf{Configuration} & \textbf{Impact under delay} \\
\midrule

SpamAssassin &
\makecell[l]{\texttt{rbl\_timeout}\\
\hspace{1em}Default: 15\,s\\
\hspace{1em}Run~3: 1.5\,s} &
Below the 2.6\,s RTT, so DNSBL lookups are abandoned before their answers return and per-message detection falls. \\
\cmidrule(l){2-3}

&
\makecell[l]{\texttt{spf\_timeout}\\
\hspace{1em}Default: 5\,s\\
\hspace{1em}Run~3: 1.5\,s} &
Below the RTT, so SPF lookups do not resolve within the scan. \\
\midrule

rspamd &
\makecell[l]{\texttt{dns.timeout}\\
\hspace{1em}Default: 1\,s\\
\hspace{1em}Tested: 1\,s} &
Sets the per-attempt DNS wait; the effective wait is
\texttt{timeout}$\times$\texttt{retransmits}. \\
\cmidrule(l){2-3}

&
\makecell[l]{\texttt{dns.retransmits}\\
\hspace{1em}Default: 5\\
\hspace{1em}Tested: 1} &
Reducing the effective DNS wait below the RTT does not drop detection: a late answer arriving while the task still runs is still used. \\
\cmidrule(l){2-3}

&
\makecell[l]{\texttt{task\_timeout}\\
\hspace{1em}Default: 8\,s\\
\hspace{1em}Run~3: 1.5\,s} &
The whole-message ceiling, below the RTT, aborts the scan before answers return; blocklist hits become \texttt{\_FAIL} timeouts and DNS volume falls. \\
\midrule

Razor &
\makecell[l]{Connect \texttt{Timeout}\\
\hspace{1em}Default: 20\,s\\
\hspace{1em}Run~3: 2\,s} &
Below the RTT, so every catalogue-server handshake aborts before its SYN-ACK returns; detection collapses to 100\% network error. \\
\cmidrule(l){2-3}

&
\makecell[l]{\texttt{can\_read}\\
\hspace{1em}Default: 15\,s\\
\hspace{1em}Unchanged} &
Governs the wait for a server's reply; a server that accepts but is slow to respond stalls the scan up to 15\,s, producing Run~2's long scan-time tail. \\
\midrule

unbound \newline (resolver) &
\makecell[l]{\texttt{discard-timeout}\\
\hspace{1em}Default: 1900\,ms\\
\hspace{1em}Configured: 9000\,ms} &
Must exceed the RTT, otherwise delayed DNSBL answers are dropped as too old and detection fails silently. \\
\cmidrule(l){2-3}

&
\makecell[l]{\texttt{infra-cache-min-rtt}\\
\hspace{1em}Default: 50\,ms\\
\hspace{1em}Configured: 4000\,ms} &
Must exceed the RTT, otherwise the resolver treats a delayed answer as a lost packet on first contact. \\

\bottomrule
\end{tabular}
\caption{Timeout settings that govern behaviour under delay, and the effect of setting them relative to the round-trip time (RTT $\approx$ 2.6\,s).}
\label{tab:config-impact}
\end{table}



\begin{table}[H]
\centering
\small
\begin{tabular}{@{}lccc@{}}
\toprule
Tool & Baseline & Run 2 & Run 3 \\
\midrule
SpamAssassin & 5.16 MB DNS & 5.42 MB DNS & 3.41 MB DNS \\
rspamd       & 18 MB total & 32 MB total & 14 MB total \\
Razor        & 3.4 MB total & 5.1 MB total & 40.2 MB total \\
\bottomrule
\end{tabular}
\caption{Summary of network traffic volume across all tools. SpamAssassin reports DNS traffic only, rspamd reports total captured traffic, and Razor reports total captured traffic at the receiver.}
\label{tab:summary-network}
\end{table}

SpamAssassin generated approximately 5 MB of DNS traffic in the baseline and Run 2, with traffic falling in Run 3 because the resolver failed after approximately 500 messages. rspamd produced the highest traffic in Run 2 (32 MB) because delayed DNS replies triggered repeated retransmissions. Razor's network volume peaks in Run~3 (40.2,MB) because repeated TCP connection failures cause \texttt{nextserver} to contact every available server, generating many SYN retransmissions.



\begin{table}[H]\centering
\footnotesize
\renewcommand{\arraystretch}{1.3}
\begin{tabular}{@{}p{5.2cm} p{7.0cm}@{}}
\toprule
\textbf{Observed (common to all three tools)} & \textbf{Reason} \\
\midrule

Detection survives delay only when timers are set above the round-trip &
Every tool assumes replies return inside a timeout set for terrestrial latency; once the round-trip exceeds it, answers arrive too late to be used \\

A timeout below the round-trip does not degrade gracefully; detection is lost or replaced by errors &
Below the RTT no reply can return in time, so a listing becomes a timeout, an error, or a failed handshake rather than a partial result \\

Scan time rises sharply under delay &
The scan blocks on network replies, so each message waits at least one round-trip, and dependent lookups wait several \\

The filter's own lookups dominate the link, far above the mail payload &
Network rules issue many DNS or catalogue queries per message, while the email itself is small by comparison \\

Delay changes lookup volume, not just speed &
A per-attempt timeout below the RTT retransmits unanswered queries, so traffic is spent on lookups that mostly time out \\

\bottomrule
\end{tabular}
\caption{Behaviour common to SpamAssassin, rspamd, and Razor under delay.}
\label{tab:results-common}
\end{table}


